\section{Conclusão}

A conclusão é a parte final do trabalho, onde se retomam os objetivos propostos na introdução e se verifica em que medida foram alcançados. Não se trata de um resumo do trabalho, mas sim de uma síntese crítica dos principais achados.

\subsection{Elementos essenciais}

\begin{enumerate}
    \item \textbf{Retomada dos objetivos}: relembre brevemente o problema de pesquisa e os objetivos (geral e específicos) definidos na introdução.

    \item \textbf{Síntese dos resultados}: apresente de forma concisa as principais descobertas, relacionando-as diretamente aos objetivos. Evite repetir dados ou gráficos já apresentados.

    \item \textbf{Contribuições do trabalho}: destaque o que seu trabalho acrescenta ao estado da arte. Qual é a relevância científica ou prática dos resultados?

    \item \textbf{Limitações}: reconheça de forma objetiva as limitações do estudo (escopo, método, dados). Isso demonstra maturidade científica.

    \item \textbf{Trabalhos futuros}: sugira desdobramentos naturais da pesquisa, indicando caminhos para investigações complementares.
\end{enumerate}

\subsection{O que evitar}

\begin{itemize}
    \item Introduzir informações, dados ou referências novas que não foram discutidos anteriormente.
    \item Repetir literalmente trechos de outros capítulos.
    \item Fazer afirmações não sustentadas pelos resultados apresentados.
    \item Usar expressões vagas como "os resultados foram satisfatórios" sem especificar em que sentido.
\end{itemize}

\subsection{Extensão recomendada}

A conclusão deve ser proporcional ao trabalho. Em geral, de uma a três páginas são suficientes para dissertações de mestrado. O importante é ser objetivo e demonstrar clareza sobre o que foi realizado e o que se aprendeu com a pesquisa.